\section{Height Gap and Curve Identification}\label{sec:interface}

The field results identify harmonic conditional means away from prescribed
strands. To identify the curve, we must determine their boundary traces on
the strand and retain those traces when the strand is stopped. This section
separates the corresponding implications. Every application below requires
the uniform versions of the preceding field and geometric estimates; a
statement for one fixed conditioning would not suffice.

\subsection{Conditional means viewed from harmonic landings}

Throughout this section $\gamma_n$ denotes the ordered \emph{dual sign
strand}, rather than the affine zero segments. Write $\mathcal F^n_t$
for the sigma-field generated by its oriented prefix up to time $t$.
It reveals adjacent signs, rather than the height ratios encoded by
affine crossing positions. Let $T$ be a stopping time for this filtration.
For a deterministic smooth test $f$ define
\[
 X_n(f)=\int_{D_n}f(z)h_n(z)\dd z,\qquad
 M^n_T(f)=\E[X_n(f)\mid\mathcal F^n_T].
\]
Thus $M^n_t(f)$ is an exact martingale. A test chosen afresh after $T$ does
not define this martingale and will not be used in the driving argument.
For a complete strand, write
$m_n(z)=\E[h_n(z)\mid\gamma_n]$.
The same notation with a superscript $T$ denotes the stopped conditional
mean. Comparisons with bounded pinned data bound these profiles uniformly.
Their subsequential limits are harmonic on the complementary open sets:
Proposition~\ref{b:cell} supplies the interior weak equation,
Proposition~\ref{b:response} gives $A>0$, and the common-wall comparison
supplies compactness.

Let $v_n$ be an interior observer vertex, with a fixed positive physical
distance from the original boundary. An independent walk from $v_n$
samples the harmonic measure
of the complete strand and the original boundary. A landing on the
positive side is understood as an oriented access; two prime ends with the
same Euclidean location need not have the same sign. We retain this side
as a label, but the two-reading comparison below does not condition on
it. That comparison conditions only on the presence of a specified
geometric incident edge and leaves the two landing signs averaged.
Each realized profile $m_n$ is still the mean conditioned on its
oriented complete strand. All boundary samples use this harmonic
sampling law, with original-boundary landings kept separate.
Lemma~\ref{g:extra} is needed to exclude additional excursions near
the incident strand.

For the global boundary-sampling statements, an experiment means a
Dobrushin domain with sign-compatible exterior heights
$+\lambda$ and $-\lambda$, where $\lambda$ is fixed as $n\to\infty$.
For any finite collection of experiments fix
$0<\lambda_0\le\min\lambda\le1/2$. Constants in estimates whose
barriers attach to original boundary may depend on $\lambda_0$.
The artificial bounded-data conditionings used for local comparison
are not additional global experiments to which boundary avoidance
is asserted.

Here is the exclusion required before using the interior records of
(G7)--(G8). Let $\xi_n$ be the walk landing and let
$d_n=\operatorname{dist}(v_n,\partial D_n)$ in lattice units.
Lemmas~\ref{g:boundary-landing} and~\ref{g:observer-landing} give
\begin{equation}\label{h:landing-exclusions}
 \lim_{\eta\downarrow0}\limsup_{n\to\infty}
 \Prob\left(\xi_n\hbox{ lies on a strand side},\
   \operatorname{dist}(\xi_n,\partial D_n\cup\{v_n\})<\eta d_n
       \right)=0.
\end{equation}
The boundary estimate is the consequence of exterior-attaching barriers
and the independent full-walk exit comparison proved there; it does
not use SLE or a matched height gap. Its stronger full-walk form
also applies to landings on an arbitrary stopped prefix.

Number the six sides of the dual hexagon at a primal vertex by
$\sigma\in\{1,\ldots,6\}$, and put
\[
 E_{n,\sigma}
 =\{\xi_n\hbox{ is a strand landing and }
                 e_\sigma^*(\xi_n)\subset\gamma_n\}.
\]
These events cover the strand landings but need not partition them.
For any finite collection of samples, first pass to a subsequence on
which all their probabilities converge. A branch with zero limiting
probability has vanishing contribution to every bounded harmonic test.
Every other branch has probability bounded below; dividing the
unnormalized exclusion error in \eqref{h:landing-exclusions} by this
lower bound gives the same exclusion under its conditional law.
Conditional on $E_{n,\sigma}$ and on $\xi_n=x$, translation by $-x$
gives exactly the law conditioned on $Z^\sigma$ from
Section~\ref{sec:geometry}. Thus the branch law is a mixture of these
translated laws. We do not choose a unique incident edge, which could
impose an additional inner event, and never condition on the landing
sign in this comparison.

\begin{lemma}[Harmonic sampling in the physical slit]\label{h:slit-sampling}
Let the original Jordan lattice domains be contained in one fixed
bounded set. Let $U_n$ be the component containing an observer $z_n$
after removing a simple dual prefix, or after cutting along a
complete simple dual crosscut. Suppose
$\operatorname{dist}(z_n,\partial U_n)\ge\rho>0$ in physical units.
Label its prime-end boundary by the original positive and negative
arcs and the two oriented strand sides; empty arcs are allowed.
There are at most four cyclic arcs before completion and at most
two in a completed component.

Kill triangular-lattice walk on its first visit to the original
boundary or to a vertex adjacent to the explored strand. Original
boundary values take priority at a vertex of both types.
Let $\xi_n$ denote this landing with its incident-side label.
There is a coupling with Brownian exit $\xi_n^{\rm B}$ from $U_n$
such that, for every $a>0$,
\[
 \sup_{U_n,z_n}
 \Prob\{|\xi_n-\xi_n^{\rm B}|>a
           \text{ or their labels differ}\}\longrightarrow0 .
 \tag{H2}
\]
The supremum includes all such prefixes, hence all stopping rules.
Repeated Euclidean boundary points are distinguished by their
incident accesses.
\end{lemma}
\begin{proof}
We first recall the required deterministic screening argument.
For a simply connected polygonal slit $U$, two complementary
prime-end arcs, and $R>0$, let $L_R$ be the union of internal
paths of diameter at most $R$ joining those arcs. The frontier
of the observer component of $U\setminus L_R$ is contained in
two balls of radius $2R$. Indeed, classify a crosscut by which
of the two arc endpoints it separates from the observer.
Two crosscuts in one class touching the observer component
must intersect, and their diameter bound then puts that entire
class in one ball. The argument uses prime ends on a slit,
not an identification of its two sides. This is the planar
argument in \cite[equation (3.55)]{SSdiscrete}.
Beurling's estimate therefore gives
\[
 \Prob_z\{\text{Brownian motion enters }L_R\}
       \le C(R/\rho)^\alpha .
\]
Apply this to each arc against its complementary arc. For a
partition into $m$ cyclic arcs, entrance into the union of
short paths joining different labels has probability at most
$Cm(R/\rho)^\alpha$, for a fixed $\alpha>0$.

The lattice walk has the required physical incidence. Every
primal edge crossed by an explored dual edge has both endpoints
observed and absorbing. Thus the interpolated walk cannot cross
the dual strand before it is killed. An interior killing vertex
has a connector of diameter $C/n$ through its own dual hexagon
to an incident strand edge, on the side bearing its sign.
If several strand edges are incident, every such access has
that same sign. At original boundary the connector is trivial
and the incoming edge specifies its access. Consecutive explored
dual edges cross primal edges belonging to a common triangle;
their absorbing endpoints form a connected set attached to the
original absorbing boundary. This set lies within $C/n$ of
the physical boundary and reaches across the annuli used below.
It supplies the connected obstacle for the discrete Beurling
estimate, including at the last explored dual vertex.

Fix $n^{-1}\ll r\ll R\ll\rho$. Couple the unrestricted
interpolated walk, with its diffusive clock, and Brownian motion
so that their distance up to a fixed time cutoff is less than
$r/4$, except with probability $o_n(1)$. The ordinary planar
invariance principle gives this coupling independently of the
slit; the triangular walk is isotropic. Both exit-time tails
at the cutoff are bounded by those for a fixed enclosing ball.
For each process use its \emph{own} first entrance into the
$2r$ boundary layer. Beurling's estimate at these two marginal
stopping times confines each remaining tail to diameter $R$,
except with total probability $C((r+n^{-1})/R)^\alpha$.
No Markov assertion about the joint coupling filtration is made.

Let $\theta$ be the first time either coupled path is within
$r$ of the physical boundary. On successful coupling both
paths are still alive, both of their $2r$ layer times have
occurred, and their positions at $\theta$ can be joined by a
segment inside $U_n$. If the exit labels differ, their confined
tails, this segment, and the lattice terminal connector give
an internal path of diameter $CR$ joining different labels.
The Brownian path has therefore entered the corresponding
mixed-access set. The preceding screening bound gives
\[
 \Prob\{\text{label mismatch or exit distance}>CR\}
 \le o_n(1)+C(r/R)^\alpha+Cm(R/\rho)^\alpha
       +\text{cutoff error}.                              \tag{H3}
\]
Take the time cutoff large, then $R$ small, then $r/R$ small,
and finally $n$ large. This proves (H2).
In particular the proof compares the walk directly with the
physical dual slit, not with an unmarked Hausdorff approximation
or a thickened domain with possibly many label changes.
\end{proof}

\begin{lemma}[Marked limits and local access]\label{h:marked-limits}
In the setting of Lemma~\ref{h:slit-sampling}, let the observers
converge, $z_n\to z$, and take a Hausdorff subsequence
of the closed removed sets consisting of the explored slit and the
original exterior. Retain the conformal maps from the disk to the
observer components, normalized by $f_n(0)=z_n$ and $f_n'(0)>0$,
and the inverse images of all boundary-label transition accesses.
After a further subsequence on which these marks converge,
bounded harmonic functions with the prescribed finite
label values converge on compact subsets to the Poisson
extension of the limiting marked arc data. Their harmonic
boundary samples converge with these marks. At a surviving
common local boundary piece of two such domains, the limit
preserves the incident-side identification.
\end{lemma}
\begin{proof}
The protected observer and the common enclosing ball make the
normalized inverse maps a nondegenerate normal family. Kernel
convergence identifies their limit with the observer component.
Extract also all transition angles on the unit circle, preserving
their cyclic order. The associated bounded step functions converge
at every angle except the finitely many limiting transitions.
Poisson integration proves compact convergence; normalized
arclength gives zero mass to those transitions. Colliding
transitions merely delete an arc of vanishing harmonic mass.
This retains the boundary marks rather than reconstructing
them from the limiting closed set.

For approach samples, use a free Brownian path independent of
all domain and field records. At fixed positive distance from
the boundary, kernel convergence and compact harmonic convergence
give convergence of its smoothed readings. The boundary-layer
tail estimate in the preceding proof lets that distance decrease
to zero after taking the mesh limit. Equivalently, the normalized
conformal maps give the limiting Brownian exit angle.

Finally suppose the two finite domains and their labels agree
on the component of a ball of radius $R$ next to a common
oriented strand. Their bounded label-Poisson extensions differ,
at a point in the concentric ball of radius $R/2$ and at distance
$r\ll R$ from that strand, by at most
$CM(r/R)^\alpha$: the difference is zero on the common local
boundary, and Beurling bounds escape from the ball.
The same bound holds if the common boundary is an original
arc. It is uniform in the domains, so it persists in marked
kernel limits. The mixed-access screening in the preceding
proof excludes conflating distinct sides in this procedure.
This proves the asserted local identification, without
claiming continuity of prime-end labels from closed sets alone.
\end{proof}

For the global experiments, the original domain approximations supply
the common bounded enclosure. Before applying these lemmas to
complete-strand samples, fix $\rho>0$ below the observer's original
boundary clearance and discard the bounded contribution from
configurations whose observer-to-strand distance is less than $\rho$
in physical units. Lemma~\ref{g:observer-landing} makes this contribution
vanish as $\rho\downarrow0$ after the mesh limit. On the retained
configurations both lemmas apply with this fixed clearance. For a
nonvanishing geometric incident-edge branch, divide the discarded
probability by its lower bound as above. This localization estimates
reading errors under the sampling law with the landing sign still
averaged; the lower comparison in (G7) remains on all its admissible
unoriented records.

For a bounded harmonic limiting profile $m$ and an independent Brownian
motion $W$ in its component, $m(W_t)$ has an almost sure limit as $W$
approaches its exit. Denote that limit by $Y$. It is a random boundary
value, not yet a deterministic height gap. We first determine its
absolute value $|Y|$ on strand landings, without conditioning on its
side label. The following elementary lemma will be applied to these
nonnegative absolute readings.

\begin{lemma}[A common lower bound for two-scale readings]
\label{h:common-kernel}
For $1\le i\le k$, let $U^i_n,V^i_n$ take values in $[-M,M]$, with
$U^i_n$ measurable with respect to a sigma-field $\mathcal A^i_n$.
Suppose that
\[
 \E|U^i_n-V^i_n|\longrightarrow0.
\]
Suppose there are probability measures $\nu_n$ on $[-M,M]$, a constant
$c>0$, and events $G^i_n\in\mathcal A^i_n$ of probability tending to one,
such that, on $G^i_n$,
\[
 \Prob(V^i_n\in B\mid\mathcal A^i_n)\ge c\,\nu_n(B)
 \quad\text{for every Borel set }B.
\]
Then every subsequence has a further subsequence on which all the
$U^i_n,V^i_n$ converge in probability to the same deterministic constant.
\end{lemma}
\begin{proof}
Pass to a subsequence on which $\nu_n$ and the laws of all $U^i_n$ converge.
Conditional integration gives
\[
 \E|U^i_n-V^i_n|
 \ge c\,\E\left[\mathbf1_{G^i_n}
          \int|U^i_n-y|\,\nu_n(\dd y)\right].
\]
The lost contribution outside $G^i_n$ is at most $2M\Prob((G^i_n)^c)$.
If $\mu_i$ and $\nu$ are the limiting laws, bounded continuity yields
$\int|x-y|\,\mu_i(\dd x)\nu(\dd y)=0$.
A product of two probability measures supported on the diagonal must
consist of the same point mass in each coordinate. The point is common
to all $i$ because $\nu$ is common. Agreement of the readings completes
the proof.
\end{proof}

The constant $c$ in this lemma need not approach one. It must, however,
stay positive as unused scale buffers grow. A different reference law for
each external configuration would not provide the stated conclusion.

\begin{lemma}[Locality of bounded conditional means]\label{h:locality}
Assume the sequentially uniform comparison and trace estimates of
Theorem~\ref{c:comparison} and Theorem~\ref{t:trace}, and the local
weak equation of Proposition~\ref{b:cell}. Consider two conditioning
systems with uniformly bounded mean profiles and the same lattice
domain and original fixed data inside a ball of radius $R$. Suppose
they share a connected oriented strand across that ball and all sign
observations inside it. For averaged observations at distance of order $r$
from that strand, supported in the concentric ball of radius $R/2$,
where $r\to\infty$ in lattice units and $r/R\to0$,
the two conditional means differ by a quantity tending to zero.

There is also an exterior version. If the two systems agree outside
a ball of radius $\rho$, their original exterior data agree, and their
unchanged strands connect that ball to radius $R$, then at distance
of order $r$ from its center, with $\rho\ll r\ll R$, the limiting
mean difference is at most $CM(\rho/r)^{1/2}$. Here $M$ bounds the
individual profiles. In both assertions use translates and dilates
of the fixed radial probability kernel from
Section~\ref{sec:trace-analysis}. A sampling ball centered at $y$
has radius $s$ less than one tenth of both $r$ and the distance from
$y$ to the boundary of the common free region. For each fixed $a>0$,
retain $s\ge ar$ while taking the mesh limit in coordinates rescaled
by $r$; the cutoff may then decrease to zero.
\end{lemma}
\begin{proof}
For the first assertion compare the two fields with common reflecting
sign walls. On shared strands their limiting mean difference has zero
Brownian-tested trace by the comparison and trace estimates. The
common original fixed data give zero differences at the corresponding
pins, so the same trace argument gives zero trace on the common
original boundary pieces. Away from these observed boundaries the
difference is harmonic by the weak equation and $A>0$. Its
absolute value on the other boundary pieces is at most $2M$.
The maximum principle therefore bounds it by $2M$ times the
probability of escaping the common ball before hitting the connected
common strand. Beurling bounds that probability by $C(r/R)^{1/2}$.
For arbitrary sequences with $r\to\infty$ and $r/R\to0$, rescale
by $r$ and first restrict to a common ball of radius $Lr$ about a
nearest strand point, with $L$ fixed. Centrality puts this ball inside
the shared region for all sufficiently large scales. Local $L^1$
compactness transfers its $CM L^{-1/2}$ bound to the fixed-cutoff
averages. First take the mesh limit, then $L\to\infty$, and finally
remove the averaging cutoff. Sequential uniformity gives the
assertion for every admissible conditioning sequence.

For the second assertion apply the identical argument in the common
exterior region, with the small changed window as the only uncontrolled
boundary. The probability of reaching that window from distance $r$
before hitting a connected common arm is at most $C(\rho/r)^{1/2}$.
The bounds persist under averaging in a sufficiently small ball about
the observation point. In both assertions the zero trace is a
conclusion of the common-wall comparison, not an equality of the
sampled wall heights.
\end{proof}

\begin{proposition}[The two absolute readings required for the height gap]
\label{h:readings}
Assume the harmonic-mean compactness and common-sign trace comparison of
Proposition~\ref{b:cell}, Theorem~\ref{c:comparison}, and
Theorem~\ref{t:trace}. Assume also the landing exclusions
\eqref{h:landing-exclusions} and the sampled-configuration estimates
of Proposition~\ref{g:sample-comparison} and Lemma~\ref{g:extra}.
Then, along any joint subsequence of finitely many experiments, after
extracting their geometric incident-edge branch probabilities as above,
the following construction gives the hypotheses of
Lemma~\ref{h:common-kernel} simultaneously for all nonvanishing branches.
Its readings take values in $[0,M]$, and their limiting variables are
the absolute values $|Y|$ of the Brownian strand-boundary samples of the
conditional mean profiles. No landing-side sign is conditioned upon.
\end{proposition}
\begin{proof}
Regard each nonvanishing event $E_{n,\sigma}$ as one member of the
finite collection; independence between these members is not needed.
If none remain, the assertion about strand samples is vacuous.
Work under each branch's conditional law, disintegrated over the landing
location and translated as above. We use the resulting $Z^\sigma$
notation, always with the landing sign still averaged. The actual
orientation remains part of each signed conditional mean.
We construct auxiliary signed inner and outer readings
$\widetilde V_n,\widetilde U_n$ which approximate the same signed
boundary sample. Their absolute values will be the readings $V_n,U_n$.
The inner absolute reading, unlike its signed auxiliary version, will
be a function of the unoriented record and so will admit the common
lower comparison from (G7).

For a target error, first use \eqref{h:landing-exclusions} to discard
strand landings within distance $\eta d_n$ of the original boundary
or observer, choosing $\eta>0$ so that their probability is small.
Recenter a retained landing at zero. Choose the lattice radii
\[
 1\ll r_{\rm in}\ll R_{\rm core}\ll R_{\rm in}\ll R_{\rm mid}
       \ll R_{\rm conn}\ll r_{\rm out}\ll n.                 \tag{H1}
\]
Require also $r_{\rm out}<\eta d_n/20$. Thus the disk
$\mathfrak B_{12R_{\rm in}}$ used by (G7) and the disks
$\mathfrak B_{4R_{\rm in}}$, $\mathfrak B_{4R_{\rm conn}}$ used by (G8)
lie inside the original domain and exclude the observer. Clearance is measurable in the outer
record, which includes the domain and observer; retaining it imposes no
additional condition on the inner strand. Write
$\mathcal A_n=\sigma(\Phi_{R_{\rm mid}})$ under this side-unconditioned
$Z^\sigma$ mixture, and let $G_n\in\mathcal A_n$ be the retained
clearance event. The diagonal choices below make $\Prob(G_n)\to1$.

\emph{Agreement of complete-profile readings.}
First take mesh limits at fixed positive physical approach distances.
In each limiting component the bounded harmonic martingale has a
boundary limit. Its values at first entrances into balls of radii
tending to zero about the eventual landing have that same limit:
these entrance times approach the exit time pathwise. They need not
be stopping times after the future-dependent center is specified.

At an interior approach point use a radial smooth average in a ball
smaller than one tenth of its distance to the strand and of the
approach radius, with the fixed probability kernel defining $P_s$ in
Section~\ref{sec:trace-analysis}. The limiting harmonic average equals
the value at its center. Impose a positive lower cutoff on the averaging radius
before taking the mesh limit, then remove it. At each fixed positive
physical approach scale the sampling point is interior almost surely,
so this cutoff loses arbitrarily little probability. The comparison
compactness and the vanishing original--filled $L^1$ error give
convergence of the conditional means in $L^1$ on every compact subset
of the limiting free regions, on a joint subsequential realization.
At a fixed radius cutoff $a>0$, the sampling supports have distance
at least $9a$ from the strand; the retained landing clearance and (H1)
also keep them a fixed positive distance from the original exterior.
Hausdorff convergence therefore places all these supports in one
compact subset $C_a$ of the limiting free regions. Since the kernels
have supremum norm at most $\|\varphi\|_\infty a^{-2}$, the discrepancy
of their readings is bounded uniformly by
$\|\varphi\|_\infty a^{-2}\|m_n-m\|_{L^1(C_a)}$, which tends to zero.
This transfers both readings to the lattice, including for kernels
selected from the records.
Denote the signed complete-profile readings at $r_{\rm in}$ and
$r_{\rm out}$ by $B_n^{\rm in}$ and $B_n^{\rm out}$, assigning zero
when their kernel cutoff fails. Such failures have vanishing probability.
A diagonal choice of the radii in (H1) makes both readings converge in
$L^1$ to $Y$ on the joint subsequential realization, and in particular
$\E|B_n^{\rm in}-B_n^{\rm out}|\to0$. These statements hold under the
side-unconditioned branch law: the bounded errors under the original
sampling law are divided only by the branch's positive probability
lower bound. The scales are chosen slowly; no convergence rate is asserted.

\emph{The reading determined by the inner record.}
Write $\theta=(\beta,w)$ for the unoriented record
$\Theta_{R_{\rm in}}$. An admissible connected incident strand through
$e_\sigma^*$ has exactly two lifts to adjacent sign observations:
choosing the sign at the landing vertex $0$ propagates the signs along
the nonbranching strand. Denote by $o_+(\theta)$ the lift for which $0$
is positive, and by $o_-(\theta)$ its sign reversal. This deterministic
choice of a lift only defines an auxiliary function of $\theta$; it
does not condition the sampled field on that sign.

For $o\in\{o_+,o_-\}$, let $m^{\rm can}_{\theta,o}$ be the conditional
mean in the centered rounded disk of radius $2R_{\rm in}$, with zero
exterior heights and precisely the strand signs in $o$. Choose the
rounding equivariantly under the six lattice rotations. These polytopes
have positive volume: assign small heights $+b$ and $-b$ to their
respective string sites and zero to the other free sites, with
$0<b<1/2$. All free edge inequalities and sign inequalities are then
strict. The volume-preserving map $h\mapsto-h$ exchanges the two
canonical polytopes and fixes the zero exterior data. Consequently
\[
 m^{\rm can}_{\theta,o_-}(z)=-m^{\rm can}_{\theta,o_+}(z)
 \quad\text{at every }z.                                      \tag{H1a}
\]

Apply (G8) at $(R_{\rm core},R_{\rm in})$ and average its conditional
bound. Except on an event of vanishing probability under $Z^\sigma$,
the walk does not leave $R_{\rm in}$ after its first entrance at
$R_{\rm core}$. Its first entrance at $r_{\rm in}$ therefore belongs
to the recorded terminal piece, and reversing that piece determines
the sampling point. Any complete-chord piece in $B(0,R_{\rm core})$
other than the incident strand $\beta_{R_{\rm in}}^\sigma$ would
give two additional crossings of the same annulus. Outside the same
type of exceptional event, the canonical conditioning with the
actual oriented lift and the complete conditioning thus agree
throughout $B(0,R_{\rm core})$.

Choose the inner radial kernel using the distance to the recorded
geometric strand and the cutoff just described; denote it by
$\kappa_\theta$. Both its center and its radius are functions only of
$\theta$. For a valid approach point and cutoff put
\[
 v_{n,o}(\theta)
   =\int\kappa_\theta(z)m^{\rm can}_{\theta,o}(z)\dd z,
 \qquad
 F_n(\theta)=|v_{n,o_+}(\theta)|.
\]
If the record has no admissible incident strand, does not contain the
required approach point, or fails the kernel cutoff, set both
$v_{n,o}$ and $F_n$ equal to zero. These fallback rules are geometric.
Equation (H1a) gives $v_{n,o_-}=-v_{n,o_+}$, including the fallback
cases. Thus
\[
 V_n:=F_n(\Theta_{R_{\rm in}})
       =|\widetilde V_n|,
 \qquad
 \widetilde V_n:=v_{n,o_{\rm actual}}(\Theta_{R_{\rm in}}).
                                                        \tag{H1b}
\]
The inner record space is countable, and the lift, kernel, and fallback
choices are deterministic. Hence $V_n$ is one bounded measurable
function of the unoriented record, the same in every experiment, and
does not use $G_n$. The absolute value is taken after testing an oriented canonical
mean; averaging the two orientations before taking the absolute value
would instead give zero and is not the construction.

Apply Lemma~\ref{h:locality} to each actual oriented conditioning,
with $r_{\rm in}/R_{\rm core}\to0$. Its sequential uniformity allows
these errors to be averaged over the actual orientations under
$Z^\sigma$. The extra-excursion estimate supplies both the common
strand neighborhood and the presence of the chosen approach point in
the recorded walk. Boundedness absorbs its exceptional events and the
cutoff failures, so
\[
 \E|B_n^{\rm in}-\widetilde V_n|\longrightarrow0.
\]
No landing sign has been conditioned on in this expectation.

\emph{The reading determined by the outer record.}
Retain the external arms and initial walk in $\Phi_{R_{\rm mid}}$.
Delete every oriented dual-edge observation whose adjacent primal
sign pair meets $B(0,R_{\rm conn})$; retain the remaining external-arm
edge observations and both vertices of each pair. Every retained sign
vertex still has an oppositely constrained neighbor.

Apply (G8) at $(R_{\rm mid},R_{\rm conn})$ and average its bound.
Except on an event of vanishing probability under $Z^\sigma$, the
unrecorded connection of the two external arms stays inside
$B(0,R_{\rm conn})$: the arms already supply two disjoint annular
crossings, and an excursion of their connection would supply two more.
Hence the modified and complete conditionings agree outside
$B(0,R_{\rm conn}+C)$, with an absolute lattice constant $C$.
All retained observations are determined by $\Phi_{R_{\rm mid}}$.
The first walk entrance at $r_{\rm out}$ is also in that record,
since it precedes the first entrance at $R_{\rm mid}$.

Use a radial kernel smaller than one tenth of $r_{\rm out}$ and
of the distance to the recorded external strands, with a positive
lower cutoff as above. Its support stays outside the changed window.
The mean under the modified conditioning, tested with this kernel,
is $\mathcal A_n$-measurable, since that record includes the oriented
external-arm observations; denote this signed reading by
$\widetilde U_n$. Set $\widetilde U_n=0$ when its required outer
clearance or kernel cutoff fails, and put $U_n=|\widetilde U_n|$.
In the common exterior component, limiting mean differences are bounded
harmonic functions with zero trace on the unchanged strands and original
boundary. The only uncontrolled boundary is the deleted window, where
the difference has absolute value at most $2M$. A connected common arm
crosses the intervening annulus. The maximum principle and Beurling's
estimate therefore bound the influence of this window at the sampling
ball by
\[
 C M(R_{\rm conn}/r_{\rm out})^{1/2}.
\]
The bounded enlargement by $C$ is negligible at these scales.
Sequentially uniform compactness and trace comparison transfer the
bound to the lattice with an error tending to zero. After including
the bounded exceptional contributions, this gives
$\E|\widetilde U_n-B_n^{\rm out}|\to0$.
The inequality $\bigl||a|-|b|\bigr|\le |a-b|$ now gives
\[
 \E|U_n-V_n|
 \le \E|\widetilde U_n-B_n^{\rm out}|
     +\E|B_n^{\rm out}-B_n^{\rm in}|
     +\E|B_n^{\rm in}-\widetilde V_n|
 \longrightarrow0.                                      \tag{H1c}
\]
All three errors are integrated under the side-unconditioned branch
law. In particular the absolute value is taken only after the signed
locality and harmonic-mean arguments.

\emph{One common lower kernel and the order of limits.}
Fix a geometric edge label $\sigma_0$. Write $F_{n,\sigma}$ for the
function just constructed in branch $\sigma$, and let $T_\sigma$ be
the lattice rotation fixing $0$ and carrying $e_\sigma^*$ to
$e_{\sigma_0}^*$. The zero exterior canonical disk, the lift making
$0$ positive, the radial kernel, and every fallback rule commute with
this rotation. A spatial change of variables therefore gives
\[
 F_{n,\sigma}(\theta)
       =F_{n,\sigma_0}(T_\sigma\theta).                  \tag{H1d}
\]
For each sufficiently large $n$, choose one admissible reference
external record $\phi_n^*$ at radius $4R_{\rm in}$ with label
$\sigma_0$ and positive hookup probability. Such a record can be
taken from any nonvanishing branch after its outer clearance and a
rotation: the clearance probability tends to one, and the discrete
record space is countable. Let $\nu_n$ be the
law on $[0,M]$ obtained by pushing
\[
 \operatorname{Law}_*(\Theta_{R_{\rm in}}\mid
             \Phi_{4R_{\rm in}}=\phi_n^*,Z^{\sigma_0})
\]
through $F_{n,\sigma_0}$.

For any fine external record $\phi$ in branch $\sigma$ with positive
hookup probability and satisfying the disk hypotheses of (G7), rotate
the experiment by $T_\sigma$ and apply the orientation-averaged (G7)
to the event
$\{F_{n,\sigma_0}(\Theta_{R_{\rm in}})\in B\}$. Equation (H1d) yields
\[
 \Prob(F_{n,\sigma}(\Theta_{R_{\rm in}})\in B
       \mid\Phi_{4R_{\rm in}}=\phi,Z^\sigma)
       \ge C^{-1}\nu_n(B)                                \tag{H1e}
\]
for every Borel set $B$. No landing-side sign is present in this
conditioning. On $G_n$, the translated domain and observer meet the
disk hypotheses of (G7) for every compatible fine record, because
that clearance is already in the outer record.

Since $R_{\rm mid}>4R_{\rm in}$, the record at $R_{\rm mid}$ is a
restriction of the record at $4R_{\rm in}$. Disintegrate also over the
original landing location in the branch mixture. Conditional
integration of (H1e) over that location and the fine record, given
$\mathcal A_n$, gives on $G_n$
\[
 \Prob(V_n\in B\mid\mathcal A_n)\ge C^{-1}\nu_n(B).
\]
The reference pushforward includes all geometric fallback records.
Neither the extra-excursion events nor successful kernel cutoffs
restrict this lower comparison. Edge-branch normalization was used
only for the averaged errors, and no factor is paid here for labels
or unused annuli.

Here is the scale order for a finite collection of limiting Brownian
boundary readings and a requested error less than $\eps>0$.
First extract the edge-branch probabilities; the finitely many
nonvanishing branches then have a common positive probability lower
bound. Choose the clearance exclusion before any physical buffers. Choose the
outer physical approach scale for both the boundary-reading error and
the retained clearance. Next choose $R_{\rm conn}/r_{\rm out}$
for exterior locality and $R_{\rm mid}/R_{\rm conn}$ for its
excursion error. Choose $R_{\rm in}\ll R_{\rm mid}$, then
$R_{\rm core}/R_{\rm in}$ for the inner excursion error, and finally
$r_{\rm in}/R_{\rm core}$ for inner locality and the inner boundary
reading. These choices are in physical units before the mesh limit, so
every retained lattice scale diverges. Hold the scales and kernel
cutoffs fixed while taking that limit, and then remove the cutoffs.
Repeat for $\eps\downarrow0$ and take a diagonal subsequence.
The comparison constant in (G7) is independent of all these unused buffers.
Lemma~\ref{h:common-kernel} now applies to $U_n,V_n$. Equation (H1c)
and the convergence of the complete readings identify their common
limiting variable with $|Y|$, not with a side-conditioned signed value.
\end{proof}

\begin{lemma}[The sign of a terminal conditional mean]\label{h:terminal-sign}
For a complete crosscut in a sign-compatible Dobrushin experiment, the
conditional mean at every positive-side primal vertex is nonnegative,
and at every negative-side primal vertex it is nonpositive. Every
limiting harmonic conditional mean has the same sign in its respective
component. If $s\in\{+1,-1\}$ is the incident-side label of a Brownian
strand-boundary sample $Y$, then $sY\ge0$.
\end{lemma}
\begin{proof}
Freeze the full vector $A$ of strand-adjacent heights in the exact
complete-strand kernel. On almost every positive-volume fiber the
remaining conditional laws are the pinned Gibbs kernels of their free
primal components, as in Section~\ref{f:section}. Every primal edge
crossed by the dual crosscut has both endpoints among the frozen string
vertices. Removing these vertices therefore disconnects the positive
and negative free graphs. A pin incident to a positive free component
is either a positive string pin or a pin on the positive original arc:
an edge from a positive free vertex to a negative pin would cross the
dual strand and would make that free vertex a frozen string vertex.
At a mark the two endpoints of the marked boundary edge are fixed, and
the crossed edge in its incident triangle has frozen string endpoints.
The exterior is fixed as well, so there is no free route around either
end of the crosscut. This also verifies the asserted incidence at the
marks.

All pins incident to a positive free component are consequently
nonnegative. Compare its conditional Gibbs kernel with the kernel on
the same component having all these pins set to zero. The monotone
single-site heat-bath coupling used in Lemma~\ref{f:bounded} orders the
two kernels, and the zero-pinned kernel has mean zero by $h\mapsto-h$.
Thus $\E[h(x)\mid\gamma_n,A]\ge0$ at each positive free vertex.
Integrating over $A$ preserves this inequality; it already holds at
positive string vertices by their observed sign. Reflection gives the
negative inequality.

These assertions are literal at primal vertices. Affine triangles
mixing the two sides lie within $O(n^{-1})$ of the dual slit. Every
compact subset of a limiting component avoids them for large $n$, so
the compact-interior convergence of the means preserves the sign
throughout that harmonic component. Its Brownian martingale and its
boundary limit preserve the same weak inequality. Lemmas
\ref{h:slit-sampling} and~\ref{h:marked-limits} identify that component
sign with the incident-side label of the sampled access.
\end{proof}

\begin{proposition}[Subsequential universal height gap]\label{h:height-gap}
Assume the uniform inputs of Proposition~\ref{h:readings} and the
trace bounds of Lemmas~\ref{g:sign-trace} and~\ref{g:upper-half}.
Every mesh subsequence has a
further subsequence and a deterministic $l\in(0,1/2]$ such that all
terminal conditional mean profiles in any prescribed finite collection
of the sign-compatible Dobrushin experiments specified above have
Brownian boundary trace $+l$ on a
positive strand side and $-l$ on a negative strand side.
\end{proposition}
\begin{proof}
Include a fixed reference Dobrushin experiment with original boundary
heights $\pm1/2$ and a fixed interior observer. Choose compact subarcs
$J_+\Subset A_+$ and $J_-\Subset A_-$ of positive length. Let
$p_n^\pm$ be the probabilities that the independent full walk, stopped
only on the original boundary, exits on their lattice approximations.
Full-domain walk convergence gives
$p_n^\pm\to\omega_D(v,J_\pm)>0$. Conditional on any completed chord,
an observer not already adjacent to it lies in one of its two
components. A full walk exiting on the opposite $J_\pm$ must cross the
chord first. The absorbing-incidence fact in
Lemma~\ref{h:slit-sampling} makes it visit a strand-adjacent vertex
before that exit. Since $J_\pm$ avoid the marks, this forced crossing
cannot be the marked boundary edge; if a crossed edge has a boundary
endpoint, its interior endpoint is already absorbing before boundary
priority can apply. An already adjacent observer has a strand landing
at time zero. Hence, for some fixed $\delta>0$ and all large $n$,
\[
 \Prob(\xi_n\hbox{ is a strand landing}\mid\gamma_n)
       \ge \min(p_n^+,p_n^-)\ge\delta.                  \tag{H1f}
\]
This is a deterministic separation argument for each chord and does
not use a barrier estimate.

Compactness supplies joint mean profiles and harmonic samples.
Proposition~\ref{h:readings} and Lemma~\ref{h:common-kernel} make
their absolute strand-boundary samples equal to one deterministic
constant $l\ge0$ on every nonvanishing geometric edge branch, shared
by all included experiments and labels. To pass back to all strand
landings, for every error event $D_n$ use the finite-union bound
\[
 \Prob(D_n,\ \xi_n\hbox{ is a strand landing})
       \le\sum_{\sigma=1}^6\Prob(D_n\cap E_{n,\sigma}).
\]
Zero-limit branches have vanishing unnormalized mass, and on every
other branch the magnitude error tends to zero. The overlaps between
the $E_{n,\sigma}$ cause no problem. Lemma~\ref{h:terminal-sign} now
gives $Y=s|Y|=sl$ on a strand access with label $s$.

For strict positivity, (H1f) and the finite cover give a reference
edge branch with positive limiting probability. Choose the
unnormalized landing-exclusion errors smaller than a fixed fraction
of that probability, and then choose the conditional (G8) errors
small. A positive limiting mass of accesses remains away from the
original boundary and observer, with only the incident strand in the
inner core and with the recorded walk approaching its accessed side.
In particular the marks, which lie on the original boundary, and the
ends of this local strand sector are outside the smaller testing scale.

Here is the access condition for applying Lemma~\ref{g:sign-trace} on
these reference sectors. In a positive completed component, freeze
all strand-adjacent heights $A$ and take the positive portion in that
lemma's Bernoulli coupling to be \emph{all} positive string pins
incident to the component. For this mixed sector take
$\lambda_*=1/2$ and $c_*=(d/2)\wedge1/2$ as in that lemma. Its pins
$B_x\in\{0,2c_*\}$ satisfy $B_x\le A_x$, and every original-arc pin
there is $+1/2\ge c_*$, the remaining comparison boundary value.
All incident pins therefore satisfy the lower-tail inequality.
By the free-graph separation in Lemma~\ref{h:terminal-sign}, the
retained graph $G$ is a union of full common unqueried components:
no free--free edge is deleted, all string walls have been frozen,
and no negative pin is incident, including at the marks. Thus
$P_{\rm u}=\varnothing$ in Lemma~\ref{g:screened-comparison}. With
fixed vertices terminal as in (SC1), every test $x\in G$ satisfies
$\delta_G(x,P_{\rm u})=\infty$, which exceeds every required
fixed-multiple access margin. The negative case follows by reflection.
This is a graph-access assertion, not a Euclidean separation from the
nearby opposite string pins.

At fixed kernel cutoff the support stays a positive fraction of the
diverging testing scale from the dual strand, so it avoids the bounded
lattice layer of frozen string vertices and lies in this unqueried
component. The exclusions and (G8) give the isolated local access at
diverging scales. Fix the oriented complete strand and the independent
walk record defining the test before sampling the auxiliary pin
transcript $(A,B)$. Given the strand, this walk record carries no
additional height information. The square cover in
Lemma~\ref{g:screened-comparison} can therefore be fixed before that
transcript, as required. Apply the pure-side trace bound of
Lemma~\ref{g:sign-trace} directly to the completed-component
original/comparison pair just described, for each actual orientation.
Its original marginal is the complete-strand kernel after averaging
the frozen heights, so this bounds the complete-profile trace itself;
the canonical open-strand pair is not used in this positivity step.
The sequentially uniform bound may then be averaged over the two
signs; it gives $sY\ge c_*>0$ on the retained positive-mass reference
accesses. This pure-access trace bound is the strict positivity input.
Since $|Y|=l$, it yields $l\ge c_*>0$.
Finally Lemma~\ref{g:upper-half}, applied to the
\emph{completed} reference profile together with
Lemma~\ref{h:terminal-sign}, gives $|Y|\le1/2$ and hence $l\le1/2$.
No upper bound for the canonical open-strand system is used.

Apply this argument to a countable collection of interior balls, with
a diagonal subsequence. If a terminal profile had a set of incorrect
signed strand-boundary values of positive harmonic measure from any
ball in that collection, its sampled value would have positive
probability of differing from $sl$ on its label-$s$ access. The finite
edge cover would detect that discrepancy on a nonvanishing branch,
contrary to the preceding conclusion. Thus the statement holds almost
surely in every component seen from that exhaustion. Original boundary
arcs retain their prescribed values by the boundary comparison,
rather than by the height-gap argument.
\end{proof}

\subsection{Passing to a stopped strand}

For an independent walk from $z$, let $\mu^z_{n,T}$ be harmonic measure on
the boundary of the component of $D_n\setminus\gamma_n[0,T]$ containing
$z$. The nonreturn input needed here is the following consequence of
Proposition~\ref{g:nonreturn}: for each compact set of protected observers,
\begin{equation}\label{h:nr}
 \lim_{s\downarrow0}\limsup_{n\to\infty}\sup_T
 \E\int_{\gamma_n[0,T]}
 \mathbf1_{\{\operatorname{dist}(x,\gamma_n(T,\mathrm{end}])<s\}}
       \,\mu^z_{n,T}(\dd x)=0.
\end{equation}
This is an annealed harmonic-mass estimate, uniformly in the stopping
rule. It does not assert that every realized stopped prefix has a
small conditional return probability.

\begin{proposition}[Stopped harmonic observable]\label{h:stopped}
Fix a subsequence furnished by Proposition~\ref{h:height-gap} and its
deterministic height gap $l$. In this proposition $n$ tends to infinity
only along that subsequence. Assume \eqref{h:nr}, the survival bound
of Lemma~\ref{g:survival}, and the sequentially uniform field
comparisons. Harmonic sampling and boundary marks are as in
Lemmas~\ref{h:slit-sampling} and~\ref{h:marked-limits}.
Let $H^l_{n,T}$ be the continuum harmonic
function in the stopped physical domain with the original exterior
values and values $+l,-l$ on the new sides. For every deterministic
smooth test $f$ supported in a protected interior ball,
\begin{equation}\label{h:stopped-observable}
 \sup_T\E\left|M^n_T(f)-\int f(z)H^l_{n,T}(z)\dd z\right|
 \longrightarrow0.
\end{equation}
The supremum is over stopping rules before a fixed neighborhood of
that ball is reached, and the assertion holds in either direction
of exploration.
\end{proposition}
\begin{proof}
It suffices to treat an arbitrary sequence $T_n$, since failure of the
supremum would select a counterexample sequence. Use one independent
full walk $S$. Write $\xi$ for its old-prefix landing, $\upsilon$
for its completed-strand landing, and set
\[
 \mathcal G_n=\sigma(\gamma_n[0,T_n],S),\qquad
 A_n=\{\xi=\upsilon\},\qquad
 q_n=\Prob(A_n\mid\mathcal G_n).
\]
For any $\mathcal G_n$-measurable $Z_n\in[0,2M]$,
\begin{equation}\label{h:reweight}
 \E Z_n\le 2M\Prob(q_n<p)+p^{-1}\E[Z_n\mathbf1_{A_n}].
\end{equation}
Indeed, on $\{q_n\ge p\}$ integrate the identity
$\E[Z_n\mathbf1_{A_n}\mid\mathcal G_n]=Z_nq_n$.
Lemma~\ref{g:survival} makes the first term arbitrarily small
with a fixed $p>0$, uniformly in $T_n$. This is the required
change of sampling law; mere absolute continuity is not used.

At physical radius $r$ let $V^{\rm old}_{n,r}$ be a small radial
average of the stopped conditional mean at the first walk entrance
into $B(\xi,r)$. Its averaging radius is less than one tenth of
the distance to the observed strand and one tenth of $r$.
Define $V^{\rm end}_{n,r}$ in the same way with the complete strand
and $\upsilon$. Positive cutoffs on the averaging radius are removed
after the mesh limit, as in Proposition~\ref{h:readings}.
Let $b(\xi)$ be $+l$ or $-l$ according to the old oriented side,
and put
\[
 Z_n=|V^{\rm old}_{n,r}-b(\xi)|
                \mathbf1_{\{\xi\notin\partial D_n\}} .
\]
This is measurable with respect to $\mathcal G_n$.
By Proposition~\ref{h:height-gap} and
Lemmas~\ref{h:slit-sampling}--\ref{h:marked-limits},
\[
 \lim_{r\downarrow0}\limsup_n
 \E\!\left[|V^{\rm end}_{n,r}-b(\upsilon)|
                \mathbf1_{\{\upsilon\notin\partial D_n\}}\right]=0.
 \tag{H4}
\]
This is the bounded harmonic martingale boundary limit in the
completed component, sampled by the lattice walk through (H2).
First discard arbitrarily small completed-observer clearances
using Lemma~\ref{g:observer-landing}; (H2) then applies with a
fixed protected observer.
The first entrance around its eventual landing need not itself
be a stopping time: these approach points tend to the exit
pathwise, which is all that this boundary-limit assertion uses.

Discard old-side landings within distance $2R$ of the original
boundary and those within distance $2R$ of the future strand.
Their total probability, denoted $e_n(R)$, satisfies
$\lim_{R\downarrow0}\limsup_n e_n(R)=0$ by
\eqref{h:landing-exclusions} and \eqref{h:nr}.
On the remaining event, and on $A_n$, the two readings use the
same walk segment, the same label, and the same averaging kernel
when $r<R/10$. The prefix extends across this common ball:
its starting point is on the distant original boundary and its
tip belongs to the distant future, or to the original boundary
if exploration is complete. The stopped and completed
conditionings have exactly the same sign observations there.
Lemma~\ref{h:locality} therefore bounds their mean difference by
$CM(r/R)^\alpha+o_n(1)$ at fixed $r,R$. Consequently
\[
 \E[Z_n\mathbf1_{A_n}]
 \le \E\!\left[|V^{\rm end}_{n,r}-b(\upsilon)|
              \mathbf1_{\{\upsilon\notin\partial D_n\}}\right]
       +CM(r/R)^\alpha+2M e_n(R)+o_n(1).                 \tag{H5}
\]
First choose the survival error in \eqref{h:reweight}, obtaining
$p>0$. Choose $R$ so that the discarded mass is small compared
with $p$ times the desired final error. Then choose $r\ll R$
to make both (H4) and the locality error equally small.
Keep these physical scales and the kernel cutoffs fixed while
taking $n\to\infty$. This order proves that the old-side boundary
reading equals its prescribed label in $L^1$, uniformly in
the chosen sequence of stopping rules.

Original boundary accesses retain their original values.
For clarity, the argument proving (GB2)--(GB3) also yields
\[
 \lim_{R\downarrow0}\limsup_n
 \Prob\{d_{D_n}(S_{\tau_{D_n}},\gamma_n)<R\}=0
 \tag{H6}
\]
in physical units: condition on the independent original-boundary
exit, apply pure-arc repulsion away from the opposite arc, and
use the two-ball bound at the two sign-change accesses.
On the complementary event the final approach has a neighborhood,
in its original-domain component, free of the completed strand.
The original-boundary comparison gives the prescribed trace there.
The pathwise coupling used for (H2), including confinement of the
terminal tails, transfers this internal-access assertion to
Brownian motion. Thus (H6) applies also on the event that the
old landing lies on the original boundary. No boundary avoidance
deduced from SLE is used.

Finally extract jointly the stopped closed sets, bounded conditional
profiles, normalized conformal maps, and their boundary marks.
The limiting stopped profile is harmonic by the weak equation.
The preceding readings give its boundary values at Brownian-almost
every marked access. Bounded harmonic uniqueness and
Lemma~\ref{h:marked-limits} identify this profile with the limit
of $H^l_{n,T_n}$. The comparison compactness and original--filled
$L^1$ estimate give convergence of the lattice conditional means
against the fixed smooth test $f$. The harmonic functions $H^l_{n,T_n}$
converge uniformly on its protected support. Integration against $f$
and boundedness therefore prove \eqref{h:stopped-observable}.
Proposition~\ref{f:averages} supplies the moments of the underlying
field averages if those averages are retained in the joint limit.

Reverse exploration reveals the same kind of sign prefixes and
satisfies the same comparisons and nonreturn estimate. Repeating
the proof gives the reverse assertion.
\end{proof}

\begin{remark}[No transfer from absolute continuity alone]
A sequence of equivalent measures can give vanishing mass to a set
whose mass stays positive under another access law. For example,
on $[0,1]$ let $\nu$ be Lebesgue measure and give $\mu_n$ density $2/n$
on $[0,1/2]$ and $2(1-1/n)$ on its complement. Then
$\mu_n([0,1/2])\to0$ while $\nu([0,1/2])=1/2$.
The survival reweighting and nonreturn estimate supply the uniform
local transfer in the preceding proof. The exact tower identity
for conditional means does not replace either estimate.
\end{remark}

\subsection{Driving convergence and the uniform curve metric}

First obtain the deterministic $l$ on a subsequence using the reference
experiment with exterior heights $\pm1/2$. Now add the experiment with
exterior heights $\pm l$ and apply Proposition~\ref{h:height-gap} to
these two experiments on that subsequence. Any further extracted
height gap must still equal $l$, since the reference absolute readings
already converge to $l$. This yields the stopped observable for
the matched experiment without assuming its existence before
$l$ was fixed.

On this height-gap subsequence we use original boundary heights
$\pm l$, matching the value on the
new sides. Map the marked domain to the upper half-plane, with the
starting point at zero and the target at infinity. If $g_t$ is its
capacity-parametrized Loewner map and $U_t$ its driver, the harmonic
observable is, with a fixed choice of side orientation,
\begin{equation}\label{h:angle}
 H^l_t(z)=l-\frac{2l}{\pi}\arg(g_t(z)-U_t).
\end{equation}
There is no boundary force-point term because the old and new heights
match.

\begin{lemma}[The observable input for driving convergence]
\label{h:driving}
For simple lattice crosscuts in marked Jordan domains, suppose that
the exact conditional-mean martingales of deterministic protected
interior averages satisfy \eqref{h:stopped-observable}, uniformly over
the stopping rules of the driving argument, with matched exterior
heights $\pm l$, where $l>0$. Suppose the comparison is to the continuum
slit-domain function \eqref{h:angle}. Then the forward chordal drivers
converge, locally uniformly in capacity time in law, to $2B$.
\end{lemma}
\begin{proof}
We specify which part of the harmonic-observable criterion is used.
Let $x_t<y_t$ be the two images of the starting point under $g_t$.
For $t>0$ introduce the bounded-driver coordinates
\[
 s=\log(y_t-x_t),\qquad
 G_s(z)=\frac{2g_t(z)-x_t-y_t}{y_t-x_t},\qquad
 \widetilde U_s=\frac{2U_t-x_t-y_t}{y_t-x_t}\in[-1,1].
\]
Let $\mathcal F_s$ be the sign-prefix filtration at the corresponding
time, equivalently the filtration of the explored dual curve after
this deterministic coordinate change. In these coordinates, as in
\cite[Sections 4.2--4.5]{SSdiscrete},
fix a finite interval $[-S,S]$ and the two deterministic interior
probes protected throughout that interval. Replace each point probe by
a radial smooth average in a sufficiently small protected disk.
The harmonic mean-value property makes the average of the continuum
observable equal to its value at the center. Thus
\eqref{h:stopped-observable} supplies the same two approximate
martingales as the point probes.

For stopping times $s_0\le s_1$ with
$s_1-s_0\le\eps^2$ and driver oscillation at most $\eps$, the two-probe
Taylor calculation gives, outside an event of arbitrarily small
prescribed probability,
\begin{align}
 \left|\E[\Delta\widetilde U+
       2\widetilde U_{s_0}\Delta s\mid\mathcal F_{s_0}]\right|
       &\le C_S\eps^3, \label{h:driver-first}\\
 \left|\E[(\Delta\widetilde U)^2-
       2(1-\widetilde U_{s_0}^{\,2})\Delta s
           \mid\mathcal F_{s_0}]\right|
       &\le C_S\eps^3. \label{h:driver-second}
\end{align}
Here $\Delta s=s_1-s_0$ and
$\Delta\widetilde U=\widetilde U_{s_1}-\widetilde U_{s_0}$.
To obtain these estimates from qualitative convergence, fix
$S,\eps$ and the exceptional probability first, and only then take
the mesh sufficiently fine. Conditional Markov inequality converts
the $L^1$ observable errors into the needed conditional errors.
No predetermined rate is required.

To apply the approximate-diffusion criterion, start the comparison
diffusion at the same bounded-driver value as the discrete process.
For a fixed comparison horizon $S$, choose the exceptional tolerance
$\eta<\delta^5/S^2$, where $\delta$ is the criterion's approximation
parameter, before choosing the mesh; Proposition~4.5 of
\cite{SSdiscrete} then gives its comparison on $[0,S-1]$.
For a limiting window $[-S,S]$, start instead at $-S'$ and let
$S'\to\infty$ after the mesh limit. The coupling argument in the
proof of Theorem~4.1 of \cite{SSdiscrete}, pp.~115--116, removes the
initial value and identifies the stationary whole-line
bounded-coordinate diffusion. Thus an arbitrary initial Jacobi law
is not substituted for the chordal limit.

The deterministic Taylor argument and the approximate-diffusion
argument of \cite[Propositions 4.2 and 4.5]{SSdiscrete} use precisely
these martingale and harmonic-observable inputs; no Gaussian
conditional law enters these steps. They identify the bounded-coordinate
diffusion with squared diffusion coefficient $2(1-u^2)$ and drift
$-2u$. The inverse coordinate change \cite[Section~4.6]{SSdiscrete} gives
the chordal driver $2B$.
The discrete harmonic-measure comparison in that source is replaced
here by the explicit continuum comparison assumed in
\eqref{h:stopped-observable}.
\end{proof}

For completeness, the normalization of this identification can also be
seen directly in any continuous Loewner subsequential limit. If the
angle at one protected point is a continuous martingale, the identity
\[
 U_t=\operatorname{Re}g_t(z)-
       \operatorname{Im}g_t(z)\cot\arg(g_t(z)-U_t)
\]
makes the driver a semimartingale after localization. Writing
$U=V+N$, where $V$ has finite variation and $N$ is a continuous local
martingale, and $Z_t=g_t(z)-U_t$, It\^o's formula gives the finite-variation
part of its angle as
\[
 -\operatorname{Im}(Z_t^{-1})\dd V_t+
 \operatorname{Im}(Z_t^{-2})
       \left(2\dd t-\tfrac12\dd[N]_t\right).
\]
Two protected points with different
$\operatorname{Re}(Z_t^{-1})$ force $\dd V_t=0$ and
$\dd[N]_t=4\dd t$. A dense set of probes provides such a pair locally.
L\'evy's characterization then gives $U=2B$.
This calculation identifies a continuous limit; the
approximate-diffusion argument is still needed to obtain driving
tightness.

\begin{proposition}[Uniform convergence of the whole curve]
\label{h:uniform-path}
Consider the matched height-gap subsequence above, with exterior heights
$\pm l$. Assume the stopped-observable conclusion of
Proposition~\ref{h:stopped} in both directions, observer avoidance
\textup{(GB4)}, harmonic-mass nonreturn \eqref{h:nr}, and the harmonic
sampling conclusion of Lemma~\ref{h:slit-sampling}. Then the interfaces
converge to chordal $\mathrm{SLE}_4$ in uniform Euclidean distance modulo
increasing reparametrization.
\end{proposition}
\begin{proof}
The argument uses more than convergence of the two endpoint drivers.
The stopped observable identifies their limiting side labels in the same
coupling. Harmonic-mass nonreturn then excludes macroscopic backtracking.
All the estimates invoked here precede this proposition.

\paragraph{Coupling the drivers and finite prefixes.}
Let $\Phi_n:\mathbb H\to D_n$ take $0$ to the starting mark and $\infty$
to the target. Normalize at a third converging boundary mark. Rad\'o's
theorem \cite[Section 1, p.~182]{Rado}, applied first in disk coordinates,
gives uniform convergence of these maps on the compactified closed
half-plane to $\Phi:\mathbb H\to D$. The hypothesis here is uniform
convergence of marked Jordan boundary parametrizations. It is stronger
than kernel convergence alone. Use the corresponding normalization for
the reverse exploration.

Write $c_n(t)$ for the forward strand in physical coordinates at
half-plane capacity $2t$, and $\widehat c_n(t)$ for its reversal with its
own capacity clock. Both clocks range over $[0,\infty]$. Indeed the final
half-edge approaches its mark non-tangentially through the polygon;
reflection at that straight boundary edge makes its image approach
infinity non-tangentially. A bounded-capacity hull has bounded height,
so the capacity tends to infinity. Put, for $t>0$,
\[
 K_{n,t}=c_n([0,t]),\qquad F_{n,t}=c_n([t,\infty]),\qquad
 L_{n,t}=\widehat c_n([0,t]).
\]
Fix a countable dense set $\mathcal Z\subset D$. Let $S_n(z)\in\{-1,1\}$
record the side of the completed crosscut containing $z$, with $+1$
for the side adjacent to the positive original arc. Assign either sign
when the classification is undefined. These are component labels, not
field signs at $z$.

Lemma~\ref{h:driving} gives both marginal driver limits $2B$.
Consequently their joint laws are tight. From every subsequence extract
further and use a Skorokhod coupling on which the two drivers converge
locally uniformly, $F_{n,t}\to F_t$ in Hausdorff distance for every
positive rational $t$, and $S_n(z)\to S(z)$ for every $z\in\mathcal Z$.
The extra records lie in countable products of compact spaces; the
finite-mesh strands can be retained with their conditional laws.
Denote the two limiting traces in physical coordinates by $c$ and
$\widehat c$. Each is marginally chordal $\mathrm{SLE}_4$: it is simple,
meets the boundary only at its endpoints, and is continuous up to the
target \cite{RohdeSchramm}. No assertion about their joint law has yet
been made.

Almost surely, for every positive rational $t$,
\begin{equation}\label{h:finite-prefixes}
 K_{n,t}\longrightarrow K_t:=c([0,t]),\qquad
 L_{n,t}\longrightarrow L_t:=\widehat c([0,t])
 \quad\hbox{in Hausdorff distance}.
\end{equation}
Here is the required deterministic consequence of driving convergence.
In half-plane coordinates, the limiting Loewner ODE is regular through
time $t$ in a neighborhood of each point outside the simple limiting
slit. The same holds at every real point other than zero. For
$\mathrm{SLE}_4$, the absolute difference between the flowed real point
and the driver, divided by two, is a Bessel process of dimension two,
which does not hit zero. First apply this to rational starting points;
real ODE order then gives the assertion for every other nonzero real
point. Continuity of the ODE in its initial point and driver excludes
prelimit hulls from small neighborhoods of all these points. For
nearby points in $\mathbb H$, the imaginary-part equation keeps the
solution in $\mathbb H$ through time $t$. Moreover, a hull at time $t$
has height at most $2\sqrt t$, and its real projection is bounded by
the driver range. Thus every subsequential limit of prefixes is
contained in the limiting slit.

Conversely, suppose a ball about an interior point of that slit were
missed by prefixes along a subsequence. Their time-$t$ conformal maps
would be holomorphic on the ball with values in $\mathbb H$.
After mapping the range to the disk, normal-family compactness and
ODE convergence off the slit give an extension of the limiting map
across the ball with values in $\mathbb H$. Its conformal inverse
would then send an interior half-plane point to the slit: approach
the slit point through the ball off the slit and use continuity of
the inverse. This is impossible. The starting point is already common
to all slits. Uniform convergence of $\Phi_n$ proves
\eqref{h:finite-prefixes} in physical coordinates. The reverse proof
is identical.

To obtain lattice stopping rules, round time $t$ up to the first
completed exploration step of capacity at least $2t$, and denote the
rule by $T_n(t)$ and its capacity time by $t'_n$. Its prefix and closed
future differ from $K_{n,t}$ and $F_{n,t}$ by at most $C/n$ in Hausdorff
distance. Also
\begin{equation}\label{h:rounded-capacity}
 t'_n\longrightarrow t\quad\hbox{almost surely}.
\end{equation}
Otherwise choose a rational $u>t$ below an overshoot occurring
infinitely often. Along those indices $K_{n,u}$ is still within one
step of $K_{n,t}$, contradicting \eqref{h:finite-prefixes}, since the
simple capacity-parametrized limit strictly extends between $t$ and
$u$. These rules are stopping times for the sign-prefix filtration.

By \textup{(GB4)}, no point of $\mathcal Z$ lies on either limiting
trace almost surely. Indeed a hit occurs in a finite prefix and,
by \eqref{h:finite-prefixes}, forces the discrete complete strand
arbitrarily close to the observer. Fatou's lemma and \textup{(GB4)}
exclude this event. All subsequent assertions can therefore be made
simultaneously for $z\in\mathcal Z$ and rational times.

\paragraph{The terminal side identifies the common limiting arc.}
For the forward limiting driver $U$, define
\[
 H_t(z)=e\left(l-\frac{2l}{\pi}
       \arg\bigl(g_t(\Phi^{-1}(z))-U_t\bigr)\right),
\]
where $e=1$ or $-1$ according to which original arc maps to the
positive real half-line. This is the harmonic observable determined
by the forward driver, with the original and new side values matched.
Equations \eqref{h:finite-prefixes}--\eqref{h:rounded-capacity} and ODE
continuity give
\begin{equation}\label{h:rounded-observable}
 H^l_{n,T_n(t)}(z)\longrightarrow H_t(z).
\end{equation}
Assign arbitrary values of absolute value at most $l$ before $z$ is
in the evaluation domain; eventually it has positive clearance from
the rounded prefix. The same statement holds in reverse.

We claim that the limiting terminal side satisfies
\begin{equation}\label{h:terminal-side}
 \E[lS(z)\mid U|_{[0,t]}]=H_t(z).
\end{equation}
Fix a small $r>0$ with $B(z,6r)\Subset D$, and take a radial smooth
averaging test $f$ of total mass one supported in $B(z,r)$.
Let $R_n$ explore to completion unless it first reaches
$\overline B(z,4r)$, when it stops, and put $P_n=T_n(t)\wedge R_n$.
Both rules protect $B(z,3r)$ for large $n$. If $\Psi$ is a bounded
continuous function of a driver on $[0,t]$, with $|\Psi|\le1$, use
at $P_n$ the multiplier $\Psi_n^*$ equal to
$\Psi(U_n|_{[0,t]})$ when capacity time $t$ has been reached, and zero
otherwise. This multiplier is $\mathcal F^n_{P_n}$-measurable. The
exact conditional means obey
\[
 \E[M^n_{R_n}(f)\Psi_n^*]=\E[M^n_{P_n}(f)\Psi_n^*].
\]
Proposition~\ref{h:stopped} replaces both conditional means by their
harmonic averages with $o_n(1)$ error, with $r$ fixed. The mean-value
property makes these averages the values at $z$. If the complete
strand avoids $\overline B(z,4r)$, the value at $R_n$ is exactly
$lS_n(z)$: the completed component has only the corresponding matched
boundary value. The other value is $H^l_{n,T_n(t)}(z)$ and the multiplier
is $\Psi(U_n|_{[0,t]})$. All harmonic values are bounded by $l$.
Consequently
\[
 \left|\E\left[(lS_n(z)-H^l_{n,T_n(t)}(z))
              \Psi(U_n|_{[0,t]})\right]\right|
 \le o_n(1)+4l\Prob\{\dist(z,\gamma_n)\le4r\}.
\]
Take the coupled limit, then $r\downarrow0$, and use \textup{(GB4)}.
Bounded continuous tests determine the conditional expectation, proving
\eqref{h:terminal-side}. Its reverse counterpart uses the same $S(z)$.
The argument applies the stopped estimate with fixed protection; it
never needs an estimate uniform as the test radius tends to zero.

For a simple limiting crosscut and $z$ off it, $H_t(z)$ tends to $+l$
or $-l$ according to the terminal side of $z$. To see this using
harmonic measure, the tail of the crosscut eventually lies in any
prescribed neighborhood of the target. Brownian motion from $z$
exits the completed side with its constant label unless it first
reaches that neighborhood. The probability of the latter event tends
to zero, since a single boundary point has zero planar Brownian hitting
probability. Along the slit the two accesses are kept distinct;
Jordan crosscuts give precisely the two prime-end sides used by the
harmonic observable.

Bounded martingale convergence in \eqref{h:terminal-side}, along integer
$t\to\infty$, now identifies $\E[lS(z)\mid U]$ with that terminal
value. A variable taking only the values $\pm l$ whose conditional
expectation is one of these two extremes equals that extreme almost
surely. Thus $S(z)$ is the side label of $c$, and also of $\widehat c$.
A Jordan crosscut is determined by its side labels on a dense set:
they determine the relative closures in $D$ of its two open sides,
whose intersection is the interior crosscut. Hence $c$ and
$\widehat c$ have the same image, traversed in opposite directions.

\paragraph{Harmonic-mass nonreturn excludes overlapping futures.}
For every positive rational $t$, we next show that almost surely
\begin{equation}\label{h:no-overlap}
 F_t\cap K_t\text{ contains no nondegenerate subarc of }K_t.
\end{equation}
Fix one $z\in\mathcal Z$, choose observer vertices $v_n\to z$, and let
$\nu_n$ be their independent-walk landing law at the rounded prefix
$T_n(t)$, conditional on the strand. Equations
\eqref{h:finite-prefixes}--\eqref{h:rounded-capacity} and
Lemma~\ref{h:slit-sampling} imply, almost surely,
\begin{equation}\label{h:landing-limit}
 \nu_n\Longrightarrow\nu,
 \qquad\nu=\text{Brownian exit law from }D\setminus K_t\text{ at }z.
\end{equation}
For clarity, the remaining Brownian domain-continuity step uses only
landing locations. Couple Brownian motions from the converging observers
by translation. Until the limiting motion approaches within $q$ of the
limiting boundary, the varying-domain paths remain inside for large
$n$. At that time their boundary distances are at most $2q$.
Beurling's estimate bounds the chance of either remaining path moving
distance $R$ before exit by $C(q/R)^\alpha$. The boundaries are connected
and have diameter bounded below, so the bound is uniform for small
$R$. Let $q/R\to0$ and then $R\to0$. This proves convergence of exit
locations and, with the slit-sampling lemma, \eqref{h:landing-limit}.

Write $F'_n$ for the closed future after $T_n(t)$, including its
current point, so that $F'_n\to F_t$. For an open set $B\Subset D$
and fixed $s>0$, weak convergence and Hausdorff convergence give
\[
 \liminf_n\nu_n\{x\in B:\dist(x,F'_n)<s\}
 \ge\nu\{x\in B:\dist(x,F_t)<s/2\}
 \ge\nu(B\cap K_t\cap F_t).
\]
For large $n$ these landings are off the original boundary. Including
the current point does not change distance to the future before
completion; completion has probability tending to zero by
\eqref{h:rounded-capacity}. Proposition~\ref{g:nonreturn}, equivalently
\eqref{h:nr} after harmonic sampling, makes the expectations of the
left-hand probabilities tend to zero as $s\downarrow0$, uniformly
in the stopping rule. Fatou's lemma and a countable exhaustion by
$B\Subset D$ give
\[
 \nu(D\cap K_t\cap F_t)=0\quad\hbox{almost surely}.
\]
Every nondegenerate subarc of the simple slit has positive harmonic
measure from $z$: a smaller subarc away from its endpoints contains
a nonempty open arc on either prime-end side, and harmonic measure
has full support on the prime-end circle. This proves
\eqref{h:no-overlap}.

\paragraph{From the common arc to ordered uniform convergence.}
Work pathwise with the preceding simultaneous conclusions. Choose
positive rational $t,p$ so large that $K_t$ and $L_p$ overlap in a
nondegenerate subarc of their common limiting arc. If, infinitely
often, the forward prefix to $t$ ended no later in strand order than
the beginning of the reverse prefix to $p$, then
$L_{n,p}\subset F_{n,t}$ on those indices. Passing to the limit would
give $L_p\subset F_t$, contradicting \eqref{h:no-overlap}. Thus the
two prefixes eventually cover the entire discrete strand. By
\eqref{h:finite-prefixes}, the complete images converge in Hausdorff
distance to the common limiting arc.

For every rational $t>0$, $F_t$ is consequently a connected compact
subset of that arc, and $K_t\cup F_t$ covers it. Hence $F_t$ contains
$c([t,\infty])$. If it extended any earlier in the arc's order,
connectedness would force a nondegenerate overlap with $K_t$, again
contrary to \eqref{h:no-overlap}. Therefore
\begin{equation}\label{h:limiting-futures}
 F_t=c([t,\infty]).
\end{equation}
These prefix and future identifications give uniform convergence of
$c_n$ to $c$ on compactified time $[0,\infty]$. Indeed, for any
$u_n\to u\in(0,\infty)$, every limit point of $c_n(u_n)$ belongs to
$K_t$ for every rational $t>u$ and to $F_t$ for every rational
$0<t<u$. Their intersection is $\{c(u)\}$. At $0$ and $\infty$ the
one-sided constraint suffices. Compactness and continuity upgrade
this sequential assertion to uniform convergence.

Every subsequence thus has a further coupled subsequence converging
in the required curve metric to chordal $\mathrm{SLE}_4$. Finally the
affine zero chord and the dual sign strand visit the same ordered
faces and differ, after increasing reparametrization, by at most
$C/n$. This transfers the conclusion to the interface of
Theorem~\ref{thm:main}.
\end{proof}

\subsection{The variance normalization}

The height-gap argument so far allows $l$ to depend on the extracted
subsequence. The conditional covariance identifies it.

\begin{proposition}[Identification of the gap]\label{h:normalization}
Assume the joint conclusions of Proposition~\ref{h:stopped},
Theorem~\ref{t:covariance}, and the field limit
Theorem~\ref{b:gff}, with the uniform moments needed to pass conditional
second moments. Then every subsequential height gap satisfies
\[
 l^2=\frac{\pi\sigma^2}{8}.
\]
\end{proposition}
\begin{proof}
Use the matched experiment with original heights $\pm l$ and the coupled
finite-prefix convergence \eqref{h:finite-prefixes}. Fix a nonzero
nonnegative smooth test $f$ supported in an interior ball. Choose a
positive rational $t$ such that a fixed neighborhood of this ball is
deterministically protected through capacity time $t$. This is possible:
its preimages under $\Phi_n$ have imaginary part bounded below, and a
hull of capacity $2t$ has height at most $2\sqrt t$. Leave a fixed
clearance so that the one-edge physical overshoot in $T_n(t)$ also
preserves protection.

Conditioning on $\mathcal F^n_{T_n(t)}$ gives exactly the sign-string
kernel of the realized oriented prefix. Whether the rule stopped there
is already decided by that prefix; no unobserved height values are
frozen. The absorbing vertices along crossed primal edges connect to
each other and to the exterior. The observed set, including the exterior
in a fixed ambient box, therefore has a connected component of uniformly
positive macroscopic diameter and satisfies the common thickness
condition in Theorem~\ref{t:covariance}. It converges in Hausdorff distance
to the exterior with $K_t$ adjoined. That theorem and
Corollary~\ref{t:random-covariance} consequently give
\begin{equation}\label{h:stopped-variance}
 \E\operatorname{Var}(X_n(f)\mid\mathcal F^n_{T_n(t)})
 \longrightarrow\sigma^2\E\langle f,G_{D\setminus K_t}f\rangle.
\end{equation}
Here brackets denote Euclidean area pairing. The Riemann-to-affine
pairing error from Theorem~\ref{b:gff} tends deterministically to zero;
the fixed-test moment bounds transfer both second moments and
conditional-mean second moments between the two pairings.

Proposition~\ref{h:stopped} approximates $M^n_{T_n(t)}(f)$ in $L^1$
by the harmonic average. ODE continuity and the protected support
identify its limit as $\int fH_t$. The exponential moments of
Proposition~\ref{f:averages}, and conditional Jensen's inequality,
make the squares of these conditional means uniformly integrable.
Thus their variances converge to $\operatorname{Var}(\int fH_t)$.
The unconditional field limit also has convergent second moments.
The law of total variance and \eqref{h:stopped-variance} yield
\begin{equation}\label{h:total-variance}
 \operatorname{Var}\left(\int fH_t\right)
 =\sigma^2\E\langle f,(G_D-G_{D\setminus K_t})f\rangle.
\end{equation}

For $0\le s\le t$ set
$Z_s(z)=g_s(\Phi^{-1}(z))-2B_s$. The imaginary part of $Z_s$ on the
test support is bounded below, since its square decreases by at most
$4s$. It\^o's formula gives
\[
 \dd\log Z_s(z)=-\frac{2}{Z_s(z)}\dd B_s.
\]
The cross variation of the harmonic observables is therefore
\begin{equation}\label{h:mean-bracket}
 \frac{16l^2}{\pi^2}
       \operatorname{Im}\frac1{Z_s(z)}
       \operatorname{Im}\frac1{Z_s(w)}\dd s.
\end{equation}
For covariance $(-\Delta)^{-1}$ the half-plane Green function is
\[
 G_{\mathbb H}(z,w)=\frac1{2\pi}
       \log\left|\frac{z-\overline w}{z-w}\right|.
\]
Differentiating its conformal pullback gives
\begin{equation}\label{h:green-loss}
 -\partial_sG_{D\setminus K_s}(z,w)=\frac2\pi
       \operatorname{Im}\frac1{Z_s(z)}
       \operatorname{Im}\frac1{Z_s(w)}.
\end{equation}
These normalizations agree with \cite[preprint version, Lemmas 2.5--2.6,
equations (2.3)--(2.4)]{SScontinuum}. Only this calculation is used;
the stronger conditional coupling in that work is not an input.

Integrate both identities against $f$ twice and over $0\le s\le t$.
The integrands are bounded on the protected support, and the Green
singularity is integrable for area pairing. No pointwise value of the
limiting field is used. Equation~\eqref{h:total-variance} becomes
\[
 \frac{16l^2}{\pi^2}\,
 \E\int_0^t\left(\int f(z)\operatorname{Im}(1/Z_s(z))\dd z\right)^2\dd s
 =
 \frac{2\sigma^2}{\pi}\,
 \E\int_0^t\left(\int f(z)\operatorname{Im}(1/Z_s(z))\dd z\right)^2\dd s.
\]
The common factor is finite and strictly positive because $f\ge0$
is nonzero and $Z_s(z)\in\mathbb H$. Cancelling it proves the formula.
\end{proof}

\begin{corollary}[Universality]\label{h:assembly}
For the model of Section~\ref{sec:model}, the constant
\[
 \lambda_*=\sigma\sqrt{\pi/8}
\]
belongs to $(0,1/2]$, is independent of the marked domain, and gives
the curve conclusion of Theorem~\ref{thm:main}. The field stiffness
is $vA=\sigma^{-2}$, and its limiting mean is the harmonic extension
of the prescribed exterior values.
\end{corollary}
\begin{proof}
Apply the reference height-gap extraction to every subsequence.
Proposition~\ref{h:normalization} gives the same value $\lambda_*$
on each extraction; its membership in $(0,1/2]$ comes from the
reference complete-crosscut bounds. Include, in each finite comparison,
the fixed experiment with exterior heights $\pm\lambda_*$.
Every subsequence then has a further subsequence to which the driving
and uniform-path conclusions apply with the same matched value.
This proves full convergence. The field assertions are the
identifications in Theorem~\ref{b:gff} and
Proposition~\ref{b:response}.
\end{proof}
